\chapter{LITERATURE REVIEW}
\pagebreak

\begin{center}
{\LARGE\textbf{LITERATURE REVIEW}}
\end{center}


This chapter reviews existing work related to IoT security, with emphasis on cyberattack classification, intrusion detection systems, honeypots, AI-based security analysis, and log visualization. The review connects these areas to the design of TrapPot and explains why the project combines Cowrie, Zeek, Random Forest classification, and the ELK Stack.

% ============================================================
\section{Thematic Analysis of IoT Security Challenges and Technical Foundations}
\label{sec:first_possible_outline}
% ============================================================

% ------------------------------------------------------------
\subsection{Challenges in Project Domain}
\label{sub:challenges_in_project_domain}
% ------------------------------------------------------------

% ------------------------------------------------------------
\subsubsection{IoT Security Vulnerabilities and Threat Landscape}
\label{subsub:iot_vulnerabilities}
% ------------------------------------------------------------

The Internet of Things faces major security risks because many devices have limited resources, weak authentication, and irregular security updates. Mahmood and colleagues reported that encryption protection was missing from 55\% of IoT devices, while authentication issues affected the device, network, and application layers \cite{mahmood2024exploring}. The growth of IoT botnets shows the practical impact of these weaknesses, as malware families such as Mirai, Gafgyt, and Mozi continue to use default passwords and exposed services to compromise devices \cite{suareztangil2023stories}. Islam and colleagues analyzed more than 11,000 IoT code examples and found that 39.58\% included real Common Vulnerabilities and Exposures instances, with 28.99\% related to denial-of-service errors that threaten network-connected IoT systems \cite{islam2023large}. These findings support TrapPot's focus on exposed remote access services and automated attack behavior.

% ------------------------------------------------------------
\subsubsection{Honeypot Deployment and Isolation Challenges}
\label{subsub:honeypot_challenges}
% ------------------------------------------------------------

Deploying honeypots for IoT security research creates both technical and operational challenges. Adaptive honeypot research has shown that fixed honeypot designs can struggle in IoT environments because they must imitate different devices while preserving realism and safety \cite{morozov2024sweet}. Network isolation is also critical because a poorly contained honeypot can become a path into production networks or a platform for attacking other systems \cite{commey2025blockchain}. Attackers may also try to detect decoy systems through behavioral analysis and timing fingerprints, so honeypots must provide believable interaction while keeping strict containment \cite{franco2021survey}. For an educational project, these constraints make a medium-interaction SSH and Telnet honeypot a practical balance between realism, safety, and implementation effort.

% ------------------------------------------------------------
\subsubsection{Machine Learning Classification Challenges}
\label{subsub:ml_challenges}
% ------------------------------------------------------------

Machine learning can help classify honeypot activity, but the approach introduces several challenges. A 2023 survey on machine-learning-based intrusion detection for IoT systems identified dataset imbalance as a major problem because benign traffic often appears in much larger quantities than rare attack classes \cite{hindy2023comprehensive}. Several IoT studies report Random Forest accuracy above 99\% for multi-class attack detection, but models can still struggle when feature sets become large or when deployment requires real-time operation on limited resources \cite{alzahrani2025lightweight}. Feature engineering therefore requires domain knowledge. Weak or irrelevant features can reduce classification quality and increase processing cost \cite{dekeersmaeker2023survey}. Model performance also depends on the quality and diversity of training data, especially when new botnet variants or zero-day behaviors differ from the examples seen during training.

% ------------------------------------------------------------
\subsubsection{Data Collection and Preprocessing Challenges}
\label{subsub:data_challenges}
% ------------------------------------------------------------

Honeypots can produce large volumes of interaction data that must be stored, cleaned, labeled, and transformed before they can support classification. Rodriguez-Gomez and colleagues noted that manual attack labeling is time-consuming, inconsistent, and difficult when attack patterns are ambiguous \cite{rodriguezgomez2025systematic}. Public IoT security datasets are useful, but they differ in collection methods, feature structures, and attack labels, which makes cross-dataset comparison difficult \cite{dekeersmaeker2023survey}. Other surveys also note that many public datasets do not fully represent the behavior captured by medium-interaction honeypots \cite{alex2023comprehensive}. For this reason, TrapPot uses public datasets for model preparation and the implemented honeypot pipeline for live operational testing.

% ------------------------------------------------------------
\subsection{Related Work}
\label{sub:related_work1}
% ------------------------------------------------------------

The following studies are closely related to TrapPot because they connect honeypots, IoT attacks, log analysis, and machine learning classification.

Lee and colleagues used Cowrie honeypots in IoT smart factory environments and applied Random Forest classification to distinguish brute-force, scanning, and malware download attacks with 96 percent accuracy and 0.1 milliseconds of inference time \cite{lee2021classification}. Their work is important to TrapPot because it validates the combination of Cowrie and Random Forest for IoT attack classification. However, TrapPot extends the project context by targeting a portable academic environment and using five classes, including command execution and normal traffic. Alajrami and Abu-Naser developed an AI-powered intrusion detection system that integrates Cowrie with Random Forest, Logistic Regression, and Isolation Forest for SSH and Telnet attack detection \cite{alajrami2025aipowered}. Their system uses a similar foundation, while TrapPot keeps the learning pipeline simpler by focusing on one trained Random Forest model and a Docker-based deployment.

Other studies validate Random Forest and ensemble learning for IoT intrusion detection without necessarily integrating them with honeypots. Saied and colleagues reported 99.99 percent accuracy for multi-class IoT botnet classification using Random Forest, outperforming several other machine learning methods \cite{saied2023comparative}. Alzahrani and colleagues developed lightweight Random Forest methods with optimized feature engineering and achieved 98 percent accuracy for IoT devices with limited resources \cite{alzahrani2025lightweight}. These studies support the algorithm choice in TrapPot, although their datasets and deployment assumptions differ from honeypot-based threat intelligence.

Research on IoT honeypot deployment also supports TrapPot's design. Saputro and colleagues used Cowrie to build medium-interaction honeypots for IoT environments and reported CPU usage below 6.3 percent, which supports the feasibility of Cowrie-based monitoring in constrained environments \cite{9359711}. Surnin and colleagues studied honeypot fingerprinting methods that attackers can use to identify decoy systems \cite{8685566}. These studies show why TrapPot must balance realism, resource use, and isolation while still providing useful data for analysis.

Systematic reviews further support the project direction. Lanz and colleagues reviewed machine learning optimization for IoT honeypots and reported that Random Forest can reach 96 percent accuracy, identifying supervised learning as an effective method for honeypot attack classification \cite{lanz2025optimizing}. Rodriguez-Gomez and colleagues emphasized that honeypot log management needs automated classification because manual analysis does not scale well \cite{rodriguezgomez2025systematic}. Franco and colleagues surveyed IoT honeypots and honeynets and confirmed that medium-interaction honeypots with Telnet and SSH emulation remain common in the field \cite{franco2021survey}. Together, these works support TrapPot's use of Cowrie, Random Forest, and structured log processing.

While existing research validates individual parts of TrapPot, the project combines them into one practical workflow: Cowrie for IoT-like SSH and Telnet deception, Zeek for network-flow evidence, Random Forest classification for five classes, and ELK dashboards for monitoring and presentation. This combination gives the project a clear educational and operational purpose within the Saudi Arabian academic context.

% ------------------------------------------------------------
\subsection{Technical Background}
\label{sub:technical_background}
% ------------------------------------------------------------

To support the system design discussed in the Methodology chapter, this section defines the main technologies and concepts used in TrapPot.

% ------------------------------------------------------------
\subsubsection{Honeypots}
\label{subsub:honeypots}
% ------------------------------------------------------------

A honeypot is a cybersecurity mechanism designed to act as a decoy system that attracts and engages attackers. By mimicking vulnerable services, it captures malicious activity without risking production assets. Honeypots are commonly categorized as low, medium, or high interaction. Medium-interaction honeypots provide a useful balance because they offer believable attacker interaction while keeping the environment controlled \cite{9268794}. They also provide threat intelligence and labeled data that can support machine-learning-based detection.

% ------------------------------------------------------------
\subsubsection{Cowrie Honeypot}
\label{subsub:cowrie}
% ------------------------------------------------------------

Cowrie is an open-source medium-interaction SSH and Telnet honeypot that emulates a UNIX-like shell environment. It records authentication attempts, executed commands, downloaded files, and session metadata \cite{9071014}. Cowrie stores events in structured formats such as JSON, which makes its logs suitable for processing, visualization, and attack classification.

% ------------------------------------------------------------
\subsubsection{SSH and Telnet Protocols}
\label{subsub:ssh_telnet}
% ------------------------------------------------------------

Secure Shell (SSH) and Telnet are network protocols used for remote command-line access. SSH provides encrypted communication, while Telnet transmits data in plaintext. Both remain important attack vectors in IoT environments because automated bots often try default or weak credentials against exposed services \cite{9558948}. The Mirai botnet compromised more than 600,000 IoT devices by exploiting weak credentials and exposed remote access services \cite{203628}.

% ------------------------------------------------------------
\subsubsection{ELK Stack}
\label{subsub:elk_stack}
% ------------------------------------------------------------

The ELK Stack consists of Elasticsearch for search and storage, Logstash for log ingestion and processing, and Kibana for visualization. The stack allows logs from different components to be processed and explored through a central interface \cite{9001068}. In TrapPot, ELK stores Cowrie, Zeek, and AI detector data and presents the results through dashboards.


% ------------------------------------------------------------
\subsubsection{Random Forest Algorithm}
\label{subsub:random_forest}
% ------------------------------------------------------------

Random Forest is an ensemble learning algorithm that combines multiple decision trees for classification. Each tree learns from a random subset of data and features, and the final class is selected through majority voting. This method reduces overfitting and performs well on structured, high-dimensional datasets. Published IoT intrusion detection studies report Random Forest accuracy in the 96 to 99 percent range while maintaining low inference latency \cite{lee2021classification,saied2023comparative}.


% ------------------------------------------------------------
\subsubsection{Zeek Network Security Monitor}
\label{subsub:Zeek}
% ------------------------------------------------------------

Zeek is an open-source network analysis framework that works as a passive network monitor. Instead of only matching signatures, Zeek extracts metadata from network traffic, such as connection duration, packet counts, byte counts, service information, and connection state. In TrapPot, Zeek observes traffic interacting with Cowrie and produces \texttt{conn.log} and \texttt{ssh.log}. These logs provide the network-level features required by the Random Forest detector.

% ------------------------------------------------------------
\subsubsection{Supervised Learning}
\label{subsub:supervised_learning}
% ------------------------------------------------------------

Supervised learning trains models using labeled datasets where inputs are paired with correct outputs. The model learns patterns during training and then applies those patterns to classify new data. TrapPot uses supervised classification because the Random Forest model learns from labeled network-traffic examples and predicts the class of new Zeek connection records.

% ------------------------------------------------------------
\subsubsection{Feature Extraction}
\label{subsub:feature_extraction}
% ------------------------------------------------------------

Feature extraction converts raw log data into numerical and categorical values that a machine learning model can use. TrapPot extracts features such as ports, protocol, service, duration, byte counts, packet counts, and connection state from Zeek logs. Well-designed features improve classification because they represent attack behavior in a compact form, while unnecessary or inconsistent features can reduce performance.


% ------------------------------------------------------------
\subsubsection{Classification Metrics}
\label{subsub:metrics}
% ------------------------------------------------------------

Model evaluation commonly uses accuracy, precision, recall, F1-score, and the confusion matrix. Accuracy measures the overall percentage of correct predictions. Precision measures how many predicted positives are correct, while recall measures how many actual positives the model detects. The F1-score balances precision and recall. A confusion matrix shows how predicted labels compare with true labels across all classes.

% ------------------------------------------------------------
\subsubsection{Attack Types}
\label{subsub:attack_types}
% ------------------------------------------------------------

TrapPot classifies five traffic classes. Brute-force attacks involve repeated username and password attempts. Scanning activity involves probing systems to identify active ports, services, and weaknesses. Malware download activity involves attempts to upload or retrieve malicious files. Command execution occurs when an attacker gains shell access and runs commands on the emulated system. Normal traffic represents baseline or benign connection behavior used by the model to separate malicious activity from ordinary network records.

% ============================================================
\section{Chronological Evolution and Progress of IoT Threat Detection Systems}
\label{sec:second_possible_outline}
% ============================================================

% ------------------------------------------------------------
\subsection{Related Work}
\label{sub:related_work2}
% ------------------------------------------------------------

IoT security became a distinct research problem in the mid-2010s as researchers documented attacks targeting Telnet and SSH services. Pa and colleagues developed IoTPOT, one of the early honeypots designed for IoT devices, and captured Telnet attacks against ARM, MIPS, and PowerPC architectures. Their work showed that Telnet scanning increased significantly from 2014 onward and confirmed that IoT devices represented a growing attack surface. Cowrie was released in 2015 as an open-source medium-interaction SSH and Telnet honeypot \cite{9071014}. Cowrie is directly relevant to TrapPot because it provides the technical foundation for the honeypot layer.

The IoT security landscape changed sharply in 2016 with the Mirai botnet, which compromised more than 600,000 IoT devices by exploiting weak default credentials on exposed remote access services. Antonakakis and colleagues documented how Mirai scanned for exposed services and used only 62 hardcoded credential pairs to gain access \cite{203628}. This study is important to TrapPot because it explains why Telnet and SSH remain central to IoT attack monitoring. Famera and colleagues later discussed Mirai variants such as Satori and Moobot, showing that the same style of weakness continues to matter \cite{famera2025mirai}.

Researchers then focused on integrating honeypots with automated classification. Saputro and colleagues deployed medium-interaction honeypots using Cowrie and showed that CPU load remained below 6.3 percent \cite{9359711}. Lee and colleagues applied Random Forest classification to IoT smart factory honeypot data and achieved 96 percent accuracy for brute-force, scanning, and malware download attacks with 0.1 milliseconds of inference latency \cite{lee2021classification}. Franco and colleagues surveyed IoT honeypots and found that medium-interaction systems prioritizing Telnet and SSH are common in the field \cite{9520645}.

Recent research from 2023 to 2025 further supports ensemble learning approaches. Saied and colleagues reported 99.99 percent accuracy for multi-class IoT botnet classification using Random Forest \cite{saied2023comparative}. Lanz and colleagues reviewed machine learning optimization for IoT honeypots and confirmed that Random Forest can reach 96 percent accuracy in this domain \cite{lanz2025optimizing}. Rodriguez-Gomez and colleagues emphasized the need for automated classification to manage large volumes of honeypot logs \cite{rodriguezgomez2025systematic}. Alajrami and Abu-Naser also developed an AI-powered intrusion detection system that integrates Cowrie with Random Forest for SSH and Telnet attack detection \cite{alajrami2025aipowered}.

TrapPot builds on this progression by adopting Cowrie's medium-interaction architecture, focusing on Telnet and SSH protocols identified by Mirai-related research, using Zeek to produce structured traffic evidence, and applying Random Forest classification supported by recent IoT detection studies. The result is a practical proof of concept for automated IoT attack monitoring that is feasible for undergraduate implementation and aligned with local cybersecurity capacity-building goals.
